\chapter*{Abstract}

This thesis addresses the critical need for reliable uncertainty quantification
in multi-label chest X-ray (CXR) classification,
where the clinical cost of errors is high.
Standard Conformal Prediction (CP) provides a baseline guarantee of marginal coverage
but fails to control for specific, clinically relevant risks such as false negatives.
To overcome this, an implementation and systematic evaluation
of two advanced conformal frameworks is presented:
Conformal Risk Control (CRC),
which controls the expected value of monotone risks like the False Negative Rate (FNR),
and Learn then Test (LTT), which offers high-probability guarantees for any risk,
including non-monotone metrics like the False Discovery Rate (FDR).

The main work involved a comparative analysis of these frameworks on large-scale CXR datasets
(CheXpert, NIH ChestX-ray14, and PadChest)
using several deep learning models (DenseNet121 and ResNet50).
Key contributions include the implementation and validation of these methods
for controlling FNR and FDR,
as well as an investigation into a clinically-weighted risk function
to prioritize critical pathologies.

Findings demonstrate that while standard CP is unsuitable
for clinical applications due to its high FNR, CRC effectively controls this risk.
Furthermore, LTT provides a more robust, high-probability guarantee for FNR
and successfully controls the non-monotone FDR, a capability beyond CRC's scope.
The results highlight a clear trade-off between the strength of the statistical guarantee
and the size of the prediction set,
offering crucial insights into the practical application of these advanced
uncertainty quantification techniques in safety-critical medical imaging tasks.
